\documentclass[11pt]{article}
\usepackage[a4paper,margin=25mm,marginparwidth=20mm,marginparsep=10pt]{geometry}
\usepackage{fontspec}
\usepackage{amsmath,amssymb,mathtools}
\usepackage{unicode-math}
\usepackage[normalem]{ulem}
\usepackage{graphicx,xcolor,tikz,tikz-cd,pdflscape,needspace}
\usetikzlibrary{calc,decorations.pathmorphing,decorations.pathreplacing,arrows.meta,shapes.geometric}
\usepackage{newunicodechar,luacode,xurl}
\usepackage[hidelinks,unicode,hypertexnames=false]{hyperref}
\setmainfont{Latin Modern Roman}
\setmathfont{Latin Modern Math}
\setmathfont{STIX Two Math}[range={"27CB,"27CD}]
% Source boldsymbol calls all request Gamma; restore its actual bold glyph.
\renewcommand\boldsymbol[1]{\symbf{#1}}
\newfontfamily\cyrfont{STIX Two Text}[Scale=MatchLowercase]
% CYRILLIC DEFINITIONS
\newunicodechar{≃}{\ensuremath{\simeq}}
\newattribute\LMattr
\newcount\LMform
\newcount\LMinApp
\LMform=1
\LMinApp=0
\newcommand\LMlanguage{fr}
\newcommand\LMcotename{}
\newcommand\LMcote[1]{\gdef\LMcotename{#1}}
\newcommand\LMnotelabel{Note marginale}
\newcommand\LMappheading{Apparat du scan}
\newcommand\LMdoubtheading{Doutes conservés dans les passages biffés}
\newcommand\LMmanuscriptlabel{Manuscrit}
\newcommand\LMcorrectedlabel{Texte}
\newcommand\LMreasonlabel{Motif}
\newcommand\LMomittedlabel{omis}
\newcommand\LMinsertedlabel{ajout intégré}
\newcommand\LMlayoutlabel{Diagramme continué sur la page en paysage}
\newcommand\LMvarref[1]{\ifmmode\else\ifinner\else\textsuperscript{\textcolor{gray}{\LMreftext{#1}}}\fi\fi}
\newcommand\LMreftext[1]{\directlua{LM.reftext("\luaescapestring{#1}")}}
\newcommand\var[4]{\ifcase\LMform #1\or #2\ifnum\LMinApp=0\LMvarref{#4}\fi\or #2\fi}
\newcommand\LMcancelledtag[1]{\ifnum\LMform=0\tag{#1}\else\tag*{#1}\fi}
\newcommand\LMadded[1]{{\color{teal}#1}}
\newcommand\deleted[1]{{\color{gray}#1}}
\newcommand\inserted[1]{{\color{teal}#1}}
\newcommand\LMarchivefooter[3]{#2}
\newcommand\LMsourceid{}
\newcommand\LMsourceMark[2]{{\def\LMsourceid{#1}#2}}
\newcommand\LMmark[1]{\directlua{LM.mark("#1",tex.inputlineno)}}
\newcommand\unsure[1]{{\color{blue}\LMmark{unsure}#1}}
\newcommand\illegible[1]{{\LMmark{illegible}\textcolor{red}{[illeg.~#1]}}}
\newcommand\LMuncertaindeletion[1]{{\color{blue}\LMmark{cancelled-anchor}\ensuremath{\dagger}}}
\newcommand\LMdoubtcopy[1]{{\color{blue}\LMmark{residual-unsure}#1}}
\newcommand\LMdoubtdata[1]{}
\newcommand\LMcancelledDoubts[1]{\ifnum\LMform=2{\color{blue}\LMmark{cancelled-anchor}\ensuremath{\dagger}}\fi}
\newcommand\scanpage[1]{%
 \par\LMflushnotes\LMflushpages\smallskip\directlua{LM.at("\luaescapestring{\LMcotename}","\luaescapestring{\detokenize{#1}}")}%
 \ifnum\LMform=2\noindent\marginpar{\footnotesize\raggedright\LMmark{scan}\textcolor{gray}{[scan #1]}}%
 \else\noindent{\footnotesize\LMmark{scan}\textcolor{gray}{[scan #1]}}\par\nobreak\fi}
\newcommand\mspage[1]{\ifnum\LMform=2\ifinner\LMqueuepage{#1}\else\marginpar{\footnotesize\raggedright\textcolor{gray}{[p.~#1]}}\fi\else{\footnotesize\textcolor{gray}{[p.~#1]}}\ \fi}
\newcommand\diagram[1]{\ifmmode\vcenter{\hbox{#1}}\else\ifinner #1\else\par\begin{center}#1\end{center}\par\LMflushnotes\LMflushpages\fi\fi}
\let\figure\relax\newcommand\figure[1]{\ifmmode\vcenter{\hbox{#1}}\else\ifinner #1\else\par\begin{center}#1\end{center}\par\LMflushnotes\LMflushpages\fi\fi}
\newcommand\LMpagefigure[2]{\ifnum\LMform=2\marginpar{\footnotesize\raggedright\textcolor{gray}{\textcircled{#2}}}\else#1\fi}
\newcommand\LMmathnotice[1]{\par\noindent{\footnotesize\color{gray}[\ensuremath{\hookrightarrow}\ #1]}\par}
\newcommand\LMlayoutnotice[1]{\par\noindent{\footnotesize\color{gray}[\LMlayoutlabel; #1]}\par}
\tolerance=1200
\emergencystretch=3em
\tracinglostchars=3
\hfuzz=0pt
\vfuzz=0pt
\raggedbottom
\setcounter{tocdepth}{1}
\setlength\parindent{1.3em}
\setlength\parskip{1.5pt plus 1pt}
\allowdisplaybreaks[2]
% LUA CODE
\makeatletter
\AtBeginDocument{\let\LMsimeq\simeq\protected\def\simeq{\ifmmode\LMsimeq\else$\LMsimeq$\fi}%
 \let\LMGamma\Gamma\protected\def\Gamma{\ifmmode\LMGamma\else$\LMGamma$\fi}}
\output\expandafter{\expandafter\unsetattribute\expandafter\LMattr\the\output}
\AtEndDocument{\clearpage\directlua{LM.summary()}}
\newif\ifLMtrial
\newcount\LMnotefresh
\newcommand\LMcurrentnoteref{}
\newcommand\LMgetnoteref{\directlua{LM.noteref("\luaescapestring{\LMsourceid}")}}
\newsavebox\LMmbox
\newdimen\LMmw
\newdimen\LMmh
\newdimen\LMnotewidth
\newcommand\LMmarginstyle{\footnotesize\rightskip=0pt plus 2.5em\relax\parfillskip=0pt plus 1fil\relax\hyphenpenalty=50\relax\def\margin##1{{[##1]}}}
\newcommand\LMtrial[2]{%
 \LMtrialtrue\directlua{LM.measure=true LM.maxw=0}%
 \setbox\LMmbox\vbox{\hsize=#1\relax\@parboxrestore\LMmarginstyle #2\par}%
 \LMtrialfalse\directlua{LM.measure=false tex.setdimen("LMmw",LM.maxw)}%
 \LMmh=\dimexpr\ht\LMmbox+\dp\LMmbox\relax}
\newcommand\LMinblock[1]{%
 \directlua{LM.blocknotes=LM.blocknotes+1}%
 \par\smallskip\noindent{\color{gray}\rule{.20\linewidth}{.35pt}}\par\nobreak
 \begingroup\LMmarginstyle\parindent=0pt\@beginparpenalty=10000\relax
 \noindent\textit{\LMnotelabel.}\enspace #1\par\endgroup\smallskip}
\newcommand\LMmargin[1]{\begingroup\LMgetnoteref
 \LMtrial{\marginparwidth}{#1}%
 \ifdim\LMmw>\marginparwidth\LMnotemark{\LMcurrentnoteref}\LMinblock{\textsuperscript{m\LMcurrentnoteref}\ #1}%
 \else\ifdim\LMmh>33pt\LMnotemark{\LMcurrentnoteref}\LMinblock{\textsuperscript{m\LMcurrentnoteref}\ #1}%
 \else\directlua{LM.shortnotes=LM.shortnotes+1}\marginpar{\LMmarginstyle #1}%
 \fi\fi\endgroup}
% This source-local wrapper retains both consecutive short notes on separate lines.
\newcommand\LMpairedmargins[1]{\begingroup
 \marginpar{\LMmarginstyle
  \def\margin##1{\LMgetnoteref\directlua{LM.shortnotes=LM.shortnotes+1}\noindent##1\par}#1}%
 \endgroup}
\newcounter{LMnoteref}
\newcommand\LMpendingnotes{}
\newcommand\LMpendingpages{}
\newcommand\LMqueuepage[1]{\ifLMtrial\else\directlua{LM.pagequeue("\luaescapestring{\LMsourceid}")}\ifnum\LMnotefresh=1\g@addto@macro\LMpendingpages{\marginpar{\footnotesize\raggedright\textcolor{gray}{[p.~#1]}}}\fi\fi}
\newcommand\LMflushpages{\ifinner\else\LMpendingpages\global\let\LMpendingpages\@empty\fi}
\newcommand\LMnotemark[1]{\ifmmode{}^{\scriptscriptstyle\text{\color{gray}m#1}}\else\textsuperscript{\color{gray}m#1}\fi}
\newcommand\LMqueuenote[1]{%
 \LMgetnoteref\LMnotemark{\LMcurrentnoteref}%
 \ifnum\LMnotefresh=1\begingroup\edef\LMnext{\endgroup\noexpand\g@addto@macro\noexpand\LMpendingnotes{\noexpand\LMqueuednote{\LMcurrentnoteref}{\the\LMattr}{\unexpanded{#1}}}}\LMnext\fi}
\newcommand\LMqueuednote[3]{\begingroup\LMattr=#2\relax\ifnum#2>0\color{blue}\fi\LMinblock{\textsuperscript{m#1}\ #3}\endgroup}
\newcommand\LMflushnotes{\ifinner\else\begingroup\let\LMnextnotes\LMpendingnotes\global\let\LMpendingnotes\@empty\LMnextnotes\endgroup\fi}
\newcommand\margin[1]{\ifmmode\LMqueuenote{\(#1\)}\else\ifinner\LMqueuenote{#1}\else\LMmargin{#1}\fi\fi}
\AddToHook{env/equation*/after}{\LMflushnotes\LMflushpages}
\AddToHook{env/equation/after}{\LMflushnotes\LMflushpages}
\AddToHook{env/align*/after}{\LMflushnotes\LMflushpages}
\AddToHook{env/align/after}{\LMflushnotes\LMflushpages}
\AddToHook{env/gather*/after}{\LMflushnotes\LMflushpages}
\AddToHook{env/center/after}{\LMflushnotes\LMflushpages}
\makeatother
\newenvironment{LMapparatus}{\par\smallskip\begingroup\LMinApp=1\footnotesize\Needspace{4\baselineskip}\csname @beginparpenalty\endcsname=10000\relax\noindent\textcolor{gray}{\rule{.26\linewidth}{.35pt}}\par\nobreak\noindent\textbf{\LMappheading}\par\nobreak\smallskip}{\par\endgroup}
\newcommand\LMentry[6]{\par\noindent{\color{gray}\texttt{#1}}\enspace\textit{#2.}\enspace\textit{\LMmanuscriptlabel:}\enspace #3\par\noindent\textit{\LMcorrectedlabel:}\enspace #4\enspace\textit{\LMreasonlabel:}\enspace #5\par\smallskip}
\newcommand\LMcarryover[2]{\par\noindent{\scriptsize\color{gray}\texttt{#1}}\enspace #2\par}
